Let $\lbrace 1, 2, 3, 4 \rbrace$ be a set of abstract symbols on which the
associative binary operation $*$ is defined by the following operation table
(associative means $(a * b) * c = a * (b * c)$):

$$
\begin{array}{|c|c|c|c|c|}
\hline
\* & 1 & 2 & 3 & 4 \\ \hline
 1 & 1 & 2 & 3 & 4 \\ \hline
 2 & 2 & 1 & 4 & 3 \\ \hline
 3 & 3 & 4 & 1 & 2 \\ \hline
 4 & 4 & 3 & 2 & 1 \\ \hline
\end{array}
$$

If the operation $*$ is represented by juxtaposition, e.g., $2 * 3$ is written
as $23$ etc., then it is easy to see from the table that of the four possible
"words" of length two that can be formed using only 2 and 3, i.e., 22, 23, 32
and 33, exactly two, 22 and 33, are equal to 1. Find a formula for the number
$A(n)$ of words of length $n$, formed by using only 2 and 3, that equal 1. From
the table and the example just given for words of length two, it is clear that
$A(1) = 0$ and $A(2) = 2$. Use the formula to find $A(12)$.
